\selectlanguage{american}%

\section{Lower Bounds\label{sec:Lower-Bounds}}

In this chapter, we have discussed a few cutting plane methods. In
particular, we showed that $O(\min(n,\sqrt{\frac{L}{\mu}})\log(\frac{1}{\epsilon}))$
gradient computations suffice. We conclude this section by showing
that $\min(n,\sqrt{\frac{L}{\mu}})$ is in fact optimal among ``gradient-based''
methods. For convenience the reader may consider $\mu=1$ and $L=\kappa$.
\begin{thm}
Fix any $L\geq\mu>0$. Consider the function
\begin{equation}
f(x)=-\frac{L-\mu}{4}x_{1}+\frac{L-\mu}{8}\sum^{n-1}_{i=1}(x_{i}-x_{i+1})^{2}+\frac{L-\mu}{8}x^{2}_{1}+\frac{\mu}{2}\sum^{n}_{i=1}x^{2}_{i}+\frac{\sqrt{L\mu}-\mu}{4}x^{2}_{n}\label{eq:worst_function}
\end{equation}

which satisfies $\mu\cdot I\preceq\nabla^{2}f(x)\preceq L\cdot I$.
Assume that our algorithm satisfies $x^{(k)}\in\mathrm{span}(x^{(0)},\nabla f(x^{(0)}),\cdots,\nabla f(x^{(k-1)}))$
with the initial point $x^{(0)}=0$. Then, for $k<n$,
\[
f(x^{(k)})-\min_{x}f(x)\geq(\frac{\mu}{L})^{3/2}(\frac{\sqrt{L/\mu}-1}{\sqrt{L/\mu}+1})^{2k}(f(x^{(0)})-\min_{x}f(x)).
\]
\end{thm}

\begin{proof}
First, we check the strong convexity and smoothness. Note that $f(x)=-x_{1}+\frac{1}{2}x^{\top}Ax$
for some matrix $A$. Hence, $\nabla^{2}f(x)=A$. Hence, we have
\[
\theta^{\top}\nabla^{2}f(x)\theta=\frac{L-\mu}{4}\sum^{n-1}_{i=1}(\theta_{i}-\theta_{i+1})^{2}+\frac{L-\mu}{4}\theta^{2}_{1}+\mu\sum^{n}_{i=1}\theta^{2}_{i}+\frac{\sqrt{L\mu}-\mu}{2}\theta^{2}_{n}
\]
For the upper bound $\nabla^{2}f(x)\preceq L\cdot I$, we note that
\begin{align*}
\theta^{\top}\nabla^{2}f(x)\theta & \leq\frac{L-\mu}{4}\sum^{n-1}_{i=1}(2\theta^{2}_{i}+2\theta^{2}_{i+1})+\frac{L-\mu}{4}\theta^{2}_{1}+\mu\sum^{n}_{i=1}\theta^{2}_{i}+(\frac{\sqrt{L\mu}-\mu}{2})\theta^{2}_{n}\\
 & \leq(L-\mu)\sum^{n}_{i=1}\theta^{2}_{i}+\mu\sum^{n}_{i=1}\theta^{2}_{i}\\
 & =L\sum^{n}_{i=1}\theta^{2}_{i}.
\end{align*}
For the lower bound $\nabla^{2}f(x)\succeq\mu\cdot I$, we note that
\[
\theta^{\top}\nabla^{2}f(x)\theta\geq\mu\sum^{n}_{i=1}\theta^{2}_{i}.
\]

To lower bound the error, we note that the gradient at $x^{(0)}$
is of the form $(?,0,0,0,\cdots)$ and hence by the assumption $x^{(1)}=(?,0,0,0,\cdots)$
and $\nabla f(x^{(1)})=(?,?,0,0,\cdots)$. By induction, only the
first $k$ coordinates of $x^{(k)}$ are non-zero.

Now, we compute the minimizer of $f(x)$. Let $x^{*}$ be the minimizer
of $f(x)$. By the optimality condition, we have that
\begin{align*}
-\frac{L-\mu}{4}+\frac{L-\mu}{4}(x^{*}_{1}-x^{*}_{2})+\frac{L-\mu}{4}x^{*}_{1}+\mu x^{*}_{1} & =0,\\
\frac{L-\mu}{4}(x^{*}_{i}-x^{*}_{i-1})+\frac{L-\mu}{4}(x^{*}_{i}-x^{*}_{i+1})+\mu x^{*}_{i} & =0,\text{ for }i\in\{2,3,\cdots,n-1\}\\
\frac{L-\mu}{4}(x^{*}_{n}-x^{*}_{n-1})+\frac{\sqrt{L\mu}+\mu}{2}x^{*}_{n} & =0.
\end{align*}
By a direct substitution, we have that $x^{*}_{i}=(\frac{\sqrt{L/\mu}-1}{\sqrt{L/\mu}+1})^{i}$
is a solution of the above equation. Now, we note that
\begin{align*}
\|x^{(k)}-x^{*}\|^{2}_{2} & \geq\sum^{n}_{i=k+1}(\frac{\sqrt{L/\mu}-1}{\sqrt{L/\mu}+1})^{2i}\geq(\frac{\sqrt{L/\mu}-1}{\sqrt{L/\mu}+1})^{2(k+1)}
\end{align*}
and that
\[
\|x^{(0)}-x^{*}\|^{2}_{2}\leq\sum^{\infty}_{i=1}(\frac{\sqrt{L/\mu}-1}{\sqrt{L/\mu}+1})^{2i}=(\frac{\sqrt{L/\mu}-1}{\sqrt{L/\mu}+1})^{2}\frac{\sqrt{L/\mu}+1}{2}.
\]
Now, by smoothness and by the strong convexity of $f$, we have
\[
\frac{f(x^{(k)})-f(x^{*})}{f(x^{(0)})-f(x^{*})}\geq\frac{\mu}{L}\cdot\frac{\|x^{(k)}-x^{*}\|^{2}_{2}}{\|x^{(0)}-x^{*}\|^{2}_{2}}\geq(\frac{\mu}{L})^{3/2}(\frac{\sqrt{L/\mu}-1}{\sqrt{L/\mu}+1})^{2k}.
\]
\end{proof}
Note that this ``worst'' function naturally appears in many problems.
So, it is a problem we need to address. In some sense, the proof points
out a common issue of any algorithm which only uses gradient information.
Given any convex function, we construct the dependence graph $G$
on the set of variables $x_{i}$ by connecting $x_{i}$ to $x_{j}$
if $\nabla f(x)_{i}$ depends on $x_{j}$ or $\nabla f(x)_{j}$ depends
on $x_{i}$ (given all other variables). Note that the dependence
graph $G$ of the worst function is simply an $n$-vertex path, whose
diameter is $n-1$. Also, note that gradient descent can only transmit
information from one vertex to another in each iteration. Therefore,
it takes at least $\Omega(\mathrm{diameter})$ time to solve the problem
unless we know the solution is sparse (when $L/\mu$ is small). However,
we note that this is not a lower bound for all algorithms.

The problem (\ref{eq:worst_function}) is a quadratic whose optimality
conditions form a Laplacian linear system, and it can be solved in
nearly linear time using spectral graph theory. \selectlanguage{english}%

